\section*{الفصل \ref{ch:g12:synthesis-strategy} --- استراتيجية التصنيع والكيمياء الخضراء}

\begin{solution}{exo:g12:synthesis-strategy:1}
إماهة الإيثين: $46.0/(28.0 + 18.0) = \qty{100}{\percent}$. التخمير:
$2 \times 46.0/180.0 = \qty{51.1}{\percent}$.
\end{solution}

\begin{solution}{exo:g12:synthesis-strategy:2}
المجموعتان كلتاهما، \ce{Z} على \omterm{def:g11:functional-groups:amine}{أمين} الغليسين \omterm{def:g11:functional-groups:carboxylic-acid}{وإستر} الميثيل على حمض
الألانين، تُوضعان في الخطوة 1 وتُنزعان في الخطوة 3.
\end{solution}

\begin{solution}{exo:g12:synthesis-strategy:3}
المبدأ 5: استعمل مذيبات وظروف تفاعل أكثر أمانًا.
\end{solution}

\begin{solution}{exo:g12:synthesis-strategy:4}
فائض الإيثانول يجعل $Q < K$، فيتقدم التفاعل أبعد في
الاتجاه المباشر. أما \omterm{def:g12:catalysis:catalyst}{الحفّاز} فيسرّع الاتجاهين بالقدر نفسه: فتُبلغ
\omterm{def:g10:reaction-progress-table:system}{الحالة النهائية} نفسها، لكن أسرع فحسب.
\end{solution}

\begin{solution}{exo:g12:synthesis-strategy:5}
نعم: لا يُختزل إلا \omterm{def:g11:functional-groups:carbonyl}{الكيتون}، إلى مجموعة \omterm{def:g11:functional-groups:alcohol}{كحول} ثانوي؛
\omterm{def:g11:functional-groups:carboxylic-acid}{والإستر} لا يتغير.
\end{solution}

\begin{solution}{exo:g12:synthesis-strategy:6}
لمزيد من \omterm{def:g11:functional-groups:carboxylic-acid}{الإستر}: فائض من أحد المتفاعلين، أو نزع الماء (مصيدة الماء)
حالما يتكوّن. وللسرعة: \omterm{def:g10:synthesis-yield:reflux}{التسخين بالارتداد}، وإضافة حفّاز حمضي.
\end{solution}

\begin{solution}{exo:g12:synthesis-strategy:7}
مع \ce{NaBH4}:
\[ \ce{CH3-CH(OH)-CH2-CH2-CO-O-CH3} . \]
ومع \ce{LiAlH4}:
\[ \ce{CH3-CH(OH)-CH2-CH2-CH2OH} \]
(بنتان-1,4-ثنائي ول)، إضافة إلى الميثانول من الجزء الآخر \omterm{def:g11:functional-groups:carboxylic-acid}{للإستر}.
\end{solution}

\begin{solution}{exo:g12:synthesis-strategy:8}
الإضافة (\qty{100}{\percent})، والأسترة
(\qty{83}{\percent}) وطريق الخطوات الثلاث إلى الإيبوبروفين
(\qty{77}{\percent}). $77.4/40.0 = 1.9$: أي أفضل بنحو مرتين.
\end{solution}

\begin{solution}{exo:g12:synthesis-strategy:9}
$180.0/(138.0 + 102.0) = \qty{75.0}{\percent}$. وحمض الإيثانويك:
$60.0/180.0 = \qty{0.333}{kg}$ لكل كيلوغرام من الأسبرين.
\end{solution}

\begin{solution}{exo:g12:synthesis-strategy:10}
يسرّع التسخين التفاعل؛ ويعيد المكثف الأبخرة
إلى الدورق، فلا يضيع \omterm{def:g7:chemical-reactions:reactant}{متفاعل} ولا \omterm{def:g6:solutions-and-solubility:solute}{مذيب}. فالتسخين يحسّن
السرعة، لا \omterm{def:g10:synthesis-yield:yield}{المردود} النهائي.
\end{solution}

\begin{solution}{exo:g12:synthesis-strategy:11}
ست خطوات مقابل ثلاث. والمبادئ: منع النفايات (1)، وجعل \omterm{def:g12:synthesis-strategy:atom-economy}{الاقتصاد الذري} أكبر ما يمكن
(2)، واستعمال الحفّازات (9)؛ وكذلك خطوات أقل، فطاقة أقل (6).
\end{solution}

\begin{solution}{exo:g12:synthesis-strategy:12}
$0.80 \times 0.70 \times 0.90 = 0.504$: \qty{50.4}{\percent}. \omterm{def:g5:raw-materials-and-recycling:raw-material}{والمادة
الأولية}: $0.10/0.504 = \qty{0.20}{mol}$.
\end{solution}

\begin{solution}{exo:g12:synthesis-strategy:13}
احمِ \omterm{def:g11:functional-groups:amine}{أمين} الألانين (\ce{Z}) وحمض الغليسين (\omterm{def:g11:functional-groups:carboxylic-acid}{إستر} الميثيل)؛
وكوّن رابطة \omterm{def:g11:functional-groups:amine}{الأميد} بين حمض الألانين الحر \omterm{def:g11:functional-groups:amine}{وأمين} الغليسين
الحر؛ ثم انزع مجموعتي الحماية كلتيهما. ثلاث مراحل، أي خمسة
تفاعلات إذا عُدّت كل حماية وكل نزع للحماية.
\end{solution}

\begin{solution}{exo:g12:synthesis-strategy:14}
يطبّق \omterm{def:g12:catalysis:catalyst}{الحفّاز} المبدأ 9. لكن \qty{250}{\celsius} تخالف
المبدأ 6، والبنزين يخالف المبدأين 3 و5، واقتصادًا ذريًا قدره
\qty{35}{\percent} يخالف المبدأ 2: $65/35 = \qty{1.9}{kg}$ من النفايات
لكل كيلوغرام من \omterm{def:g7:chemical-reactions:reactant}{الناتج}. فسمة خضراء واحدة لا تجعل العملية
خضراء.
\end{solution}

\begin{solution}{exo:g12:synthesis-strategy:15}
% ledger: crit:CO2-T, crit:CO2-p
\qty{40}{\celsius} و\qty{100}{bar} فوق النقطة الحرجة
لثنائي أكسيد الكربون، \qty{31}{\celsius} و\qty{73.8}{bar}. والضغط
أكبر من الضغط الجوي مئة مرة: فالوعاء من الفولاذ السميك. وثنائي أكسيد الكربون
ليس سامًا ولا قابلًا للاشتعال، ويعود غازًا
دون أن يترك بقايا، ويمكن إعادة استعماله؛ ولا يتسرب أي \omterm{def:g6:solutions-and-solubility:solute}{مذيب} مكلور.
\end{solution}

\begin{solution}{pb:g12:synthesis-strategy:1}
% ledger: ibu:bhc-ae-recovered
\textbf{1.} C: $10 + 4 + 1 = 15 = 13 + 2$؛ H: $14 + 6 + 2 = 22 = 18 + 4$؛
O: $3 + 1 = 4 = 2 + 2$.

\textbf{2.} حمض الإيثانويك، \ce{C2H4O2}.

\textbf{3.} النيتروجين والكلور والصوديوم: فالإيبوبروفين لا يحوي إلا الكربون
والهيدروجين والأكسجين.

\textbf{4.} الأحدث: المبدأ 9، الحفّازات بدل الكواشف
بكميات ستوكيومترية.

\textbf{5.} $13 \times 12.0 + 18 \times 1.0 + 2 \times 16.0 =
\qty{206.0}{g/mol}$.

\textbf{6.} $134.0 + 102.0 + 122.5 + 68.0 + 19.0 + 33.0 + 36.0 =
\qty{514.5}{g/mol}$.

\textbf{7.} $206.0/514.5 = \qty{40.0}{\percent}$.

\textbf{8.} $134.0 + 102.0 + 2.0 + 28.0 = \qty{266.0}{g/mol}$.

\textbf{9.} $206.0/266.0 = \qty{77.4}{\percent}$.

\textbf{10.} $(206.0 + 60.0)/266.0 = \qty{100}{\percent}$ على الورق؛ وأكثر من
\qty{99}{\percent} في الواقع.

\textbf{11.} $514.5/206.0 = \qty{2.50}{t}$.

\textbf{12.} $2.50 - 1 = \qty{1.50}{t}$.

\textbf{13.} $266.0/206.0 = \qty{1.29}{t}$.

\textbf{14.} $60.0/206.0 = \qty{0.29}{t}$.

\textbf{15.} $1.50 - 0.29 = \qty{1.21}{t}$.

\textbf{16.} $0.90^6 = 0.53$: \qty{53}{\percent}.

\textbf{17.} $0.90^3 = 0.73$: \qty{73}{\percent}.

\textbf{18.} كل خطوة تستعمل مذيبات وعمليات فصل وتنقية،
وتفقد بعض \omterm{def:g7:chemical-reactions:reactant}{الناتج}: فخطوات أقل تعني خسائر أقل من هذا النوع.

\textbf{19.} لا شيء تقريبًا: فناتجه الثانوي الوحيد يُستعمل.

\textbf{20.} يتجنب طريق الخطوات الثلاث نحو \qty{1.5}{t} من النفايات لكل
طن من الإيبوبروفين (و\qty{1.2}{t} حتى لو رُمي حمض الإيثانويك
\omterm{def:g7:chemical-reactions:reactant}{الناتج} عنه).
\end{solution}
